\begin{problem}[第34届物理决赛][磁场中的正方形线框]{83}
    如图，在磁感应强度大小为 $B$、方向竖直向上的匀强磁场中，有一均质刚性导电的正方形线框 $abcd$，线框质量为 $m$，边长为 $l$，总电阻为 $R$。线框可绕通过 $ad$ 边和 $bc$ 边中点的光滑轴 $OO'$ 转动。$P$、$Q$ 点是线框引线的两端，$OO'$ 轴和 $X$ 轴位于同一水平面内，且相互垂直。不考虑线框自感。
    \begin{subproblem}
        \begin{subsubproblem}
            求线框绕 $OO'$ 轴的转动惯量 $J$；
        \end{subsubproblem}
        \begin{subsubproblem}
            $t=0$ 时，线框静止，其所在平面与 $X$ 轴有一很小的夹角 $\theta_0$，此时给线框通以大小为 $I$ 的恒定直流电流，方向沿 $P\rightarrow a\rightarrow b\rightarrow c\rightarrow d\rightarrow Q$，求此后线框所在平面与 $X$ 轴的夹角 $\theta$、线框转动的角速度 $\dot{\theta}$ 和角加速度 $\ddot{\theta}$ 随时间变化的关系式；
        \end{subsubproblem}
        \begin{subsubproblem}
            $t=t_0>0$ 时，线框平面恰好逆时针转至水平，此时断开 $P$、$Q$ 与外电路的连接，此后线框如何运动？求 $P$、$Q$ 间电压 $V_{\mathrm{PQ}}$ 随时间变化的关系式；
        \end{subsubproblem}
        \begin{subsubproblem}
            线框做上述运动一段时间后，当其所在平面与 $X$ 轴夹角为 $\theta_{1}\ (\frac{\pi}{4}\leq\theta_{1}\leq\frac{3\pi}{4})$ 时，将 $P$、$Q$ 短路，线框再转一小角度 $\alpha$ 后停止，求 $\alpha$ 与 $\theta_{1}$ 的关系式和 $\alpha$ 的最小值。
        \end{subsubproblem}
    \end{subproblem}
\end{problem}